% Lecture notes: Week 13, Lecture 1 -- Thermal Expansion
% To hide this lecture from the website, add a line containing only:  % draft
% To set the date shown on the website, add a line like:  % posted: 2026-09-02
\documentclass[11pt]{article}
\usepackage{mech2030notes}
\begin{document}

% ##################################################################
%  WEEK 13, LECTURE 1
% ##################################################################
\lecture{13}{1}{Thermal Expansion}

\noindent
$\rightarrow$ The vast majority of engineering materials \emph{expand} when \emph{heated}.

\begin{center}
\begin{tikzpicture}
  % small cube (T1, V1)
  \begin{scope}
    \draw[thick, fill=gray!12] (0,0) rectangle (1,1);
    \draw[thick, fill=gray!25] (0,1) -- (0.35,1.3) -- (1.35,1.3) -- (1,1) -- cycle;
    \draw[thick, fill=gray!35] (1,0) -- (1.35,0.3) -- (1.35,1.3) -- (1,1) -- cycle;
    \node[below] at (0.6,-0.1) {$T_1,\ V_1$};
  \end{scope}
  % arrow
  \draw[-{Stealth[length=2.5mm]}, thick] (2.0,0.9) to[bend left=25] (4.0,0.9);
  \node[above] at (3.0,1.35) {$\Delta T > 0$};
  % large cube (T2, V2)
  \begin{scope}[shift={(4.6,-0.3)}]
    \draw[thick, fill=gray!12] (0,0) rectangle (1.6,1.6);
    \draw[thick, fill=gray!25] (0,1.6) -- (0.55,2.05) -- (2.15,2.05) -- (1.6,1.6) -- cycle;
    \draw[thick, fill=gray!35] (1.6,0) -- (2.15,0.45) -- (2.15,2.05) -- (1.6,1.6) -- cycle;
    \node[below] at (1.0,-0.1) {$T_2,\ V_2$};
  \end{scope}
\end{tikzpicture}
\end{center}

\begin{itemize}[label=$\rightarrow$]
  \item Here $T_2 > T_1$ and $V_2 > V_1$: \emph{heating} leads to \emph{volumetric
        expansion}.
  \item For \emph{slender}, \emph{beam-like} members, expansion along the long axis
        dominates.
\end{itemize}

\begin{center}
\begin{tikzpicture}
  % bar at T1
  \begin{scope}
    \lwall{0}{-0.7}{0.7}
    \draw[beam] (0,-0.15) rectangle (3,0.15);
    \node[above] at (1.5,0.15) {$T_1$};
    \draw[dim] (0,-0.6) -- (3,-0.6) node[midway, fill=white] {$L_1$};
  \end{scope}
  \draw[-{Stealth[length=2.5mm]}, thick] (3.5,0.3) to[bend left=25] (5.1,0.3);
  \node[above] at (4.3,0.7) {$\Delta T > 0$};
  % bar at T2
  \begin{scope}[shift={(5.7,0)}]
    \lwall{0}{-0.7}{0.7}
    \draw[beam] (0,-0.15) rectangle (3.8,0.15);
    \draw[thin, dashed] (3,-0.25) -- (3,0.25);
    \node[above] at (1.9,0.15) {$T_2$};
    \draw[dim] (0,-0.6) -- (3.8,-0.6) node[midway, fill=white] {$L_2$};
  \end{scope}
\end{tikzpicture}
\end{center}

\noindent
$\hookrightarrow$ For beams, the length change is
\[
  \fbox{\fbox{$\Delta L = \alpha\,\Delta T\,L_1$}} \quad \text{(thermal expansion)}
\]
and
\[
  \fbox{\fbox{$\displaystyle \epsilon_T = \frac{\Delta L}{L_1} = \alpha\,\Delta T$}}
  \quad \text{(thermal axial strain)}
\]
\begin{itemize}[label=\then]
  \item $\alpha$ is the \emph{linear coefficient of thermal expansion} (CTE).
  \begin{itemize}[label=\then]
    \item Units: $1/{}^\circ\mathrm{C}$, or $1/\mathrm{K}$, or $1/{}^\circ\mathrm{F}$.
  \end{itemize}
\end{itemize}

% ==================================================================
\section*{Example}
% ==================================================================

\noindent
\begin{minipage}[c]{0.55\textwidth}
\centering
\begin{tikzpicture}
  % cross-section
  \draw[thick, fill=gray!12] (-1.9,-0.25) rectangle (-1.4,0.25);
  \draw[dim] (-2.15,-0.25) -- (-2.15,0.25);
  \node[left] at (-2.2,0) {\small$\SI{20}{mm}$};
  \draw[dim] (-1.9,-0.55) -- (-1.4,-0.55) node[midway, below] {\small$\SI{20}{mm}$};
  % bar between walls
  \lwall{0}{-0.8}{0.8}
  \rwall{4}{-0.8}{0.8}
  \draw[beam] (0,-0.3) rectangle (4,0.3);
  \node[pin] at (0,0) {}; \node[above left] at (-0.35,0.3) {$A$};
  \node[pin] at (4,0) {}; \node[above right] at (4.35,0.3) {$B$};
  \node[below] at (2,-0.35) {$T_1 = \SI{25}{\celsius}$};
  \draw[dim] (0,1.1) -- (4,1.1) node[midway, fill=white] {$L = \SI{1}{m}$};
\end{tikzpicture}
\end{minipage}%
\hfill
\begin{minipage}[c]{0.42\textwidth}
$\rightarrow$ Square steel bar with
\[
  \left\{
  \begin{aligned}
    E &= \SI{210}{GPa} \\
    \alpha &= 12 \times 10^{-6}\ /{}^\circ\mathrm{C}
  \end{aligned}
  \right.
\]
\textbf{Q:} If the bar is heated to $T_2 = \SI{40}{\celsius}$, what are the reaction
forces?
\end{minipage}

\subsection*{FBD}

\begin{center}
\begin{tikzpicture}
  \draw[thick] (0,0) -- (4,0);
  \node[pin] at (0,0) {}; \node[pin] at (4,0) {};
  \draw[force] (-1.2,0) node[left] {$A_x$} -- (-0.08,0);
  \draw[force] (4.08,0) -- (5.2,0) node[right] {$B_x$};
\end{tikzpicture}
\end{center}

\[
  \textstyle\sum F_x = 0 \quad\rightarrow\quad \boxed{A_x + B_x = 0}
  \qquad \text{(statically indeterminate)}
\]

\noindent
$\rightarrow$ Compatibility condition:
\[
  \boxed{\delta_B = 0}
\]

\noindent
$\rightarrow$ Superposition:
\[
  \delta_B = \delta_B^{T} + \delta_B^{B_x}
\]

\begin{center}
$\delta_B \;=$\quad
\begin{tikzpicture}[baseline=(b)]
  \coordinate (b) at (0,0);
  \lwall{0}{-0.7}{0.7}
  \draw[beam] (0,-0.3) rectangle (3,0.3);
  \draw[thin, dashed] (3,0.3) -- (3.7,0.3) -- (3.7,-0.3) -- (3,-0.3);
  \node[pin] at (3,0) {}; \node[left] at (2.95,0) {$B$};
  \draw[force] (3,0.6) -- (3.7,0.6) node[above] {$\delta_B^T$};
  \node[align=center, font=\small] at (1.5,-1.0) {(thermal expansion)};
\end{tikzpicture}
\quad{\Large $+$}\quad
\begin{tikzpicture}[baseline=(b)]
  \coordinate (b) at (0,0);
  \lwall{0}{-0.7}{0.7}
  \draw[beam] (0,-0.3) rectangle (3,0.3);
  \draw[thin, dashed] (3,0.3) -- (3.7,0.3) -- (3.7,-0.3) -- (3,-0.3);
  \node[pin] at (3,0) {}; \node[left] at (2.95,0) {$B$};
  \draw[force] (3,0.6) -- (3.7,0.6) node[above] {$\delta_B^{B_x}$};
  \draw[force] (3.05,0) -- (4.3,0) node[right] {$B_x$};
  \node[align=center, font=\small] at (1.5,-1.0) {(force $B_x$)};
\end{tikzpicture}
\end{center}

\[
  \delta_B = \alpha\,\Delta T\,L + \frac{B_x L}{EA} = 0
\]
\[
  \hookrightarrow\quad
  B_x = -\frac{\alpha(\Delta T)\,L\,EA}{L}
  = \boxed{-\alpha(\Delta T)\,EA}
\]
\[
  B_x = -\left(12 \times 10^{-6}\ /{}^\circ\mathrm{C}\right)(\SI{15}{\celsius})
         (\SI{210e9}{Pa})(\SI{0.02}{m})^2
  \qquad\Rightarrow\qquad
  \fbox{\fbox{$B_x = -\SI{15.1}{kN}$}}
\]
\[
  \text{thus}\qquad
  \fbox{\fbox{$A_x = -B_x = \SI{15.1}{kN}$}}
\]

\noindent
\textbf{Final FBD} (forces in their actual directions; the bar is in compression):

\begin{center}
\begin{tikzpicture}
  \draw[thick] (0,0) -- (4,0);
  \node[pin] at (0,0) {}; \node[pin] at (4,0) {};
  \node[below] at (0,-0.05) {$A$};
  \node[below] at (4,-0.05) {$B$};
  \draw[force] (-1.2,0) -- (-0.08,0);
  \node[above] at (-0.6,0.05) {$\SI{15.1}{kN}$};
  \draw[force] (5.2,0) -- (4.08,0);
  \node[above] at (4.6,0.05) {$\SI{15.1}{kN}$};
\end{tikzpicture}
\end{center}

\end{document}
